\scanpage{49}
\margin{48}
\mspage{-3.3-}

$0\leq\alpha_i\leq m$, let
\[
D_{\alpha_0,\ldots,\alpha_n}=\det_{0\leq i,j\leq n}\bigl[X_i^{(\alpha_j)}\bigr]
\]
be the section of $L_{\alpha_0}\otimes\cdots\otimes L_{\alpha_n}$ obtained by addition and tensor product as in the determinant. It is evident that $D_{\alpha_0,\ldots,\alpha_n}$ is an \uline{invariant} section of $L_{\alpha_0}\otimes\cdots\otimes L_{\alpha_n}$. The non-vanishing of suitable $D$'s defines the open sets we are looking for.

\uline{Definition 3.2}: An \uline{$R$-partition} of $\{0,1,\ldots,n\}$ is an ordered set of subsets $S_1,S_2,\ldots,S_\nu$ of $\{0,1,\ldots,n\}$ such that

(i) $S_i\cap(S_{i-1}\cup\cdots\cup S_2\cup S_1)$ consists of exactly one integer, for $2\leq i\leq\nu$,

(ii) $\displaystyle\bigcup_{i=1}^{\nu}S_i=\{0,1,\ldots,n\}$,

\uline{Definition 3.3}: Given an $R$-partition $R=\{S_1,\ldots,S_\nu\}$, let $U_R\subset(\mathbb{P}_n)^{n+\nu+1}$ be the open subset defined by
\[
\text{(i)}\quad D_{0,1,\ldots,n}\ne0,
\]
(ii) for all $k$ between 1 and $\nu$, and for all $i\in S_k$
\[
D_{0,1,\ldots,\widehat i,\ldots,n,n+k}\ne0
\]

Not only is $U_R$ affine, but the whole structure of the action of $\mathrm{PGL}(n)$ on $U_R$ can be described explicitly:

\uline{Proposition 3.1}: Let $R=\{S_1,\ldots,S_\nu\}$ be an $R$-partition of $\{0,1,\ldots,n\}$. Let $\mathrm{PGL}(n)$ act on $\mathrm{PGL}(n)\times\mathbb{A}^{n\nu-n}$ by the product of left translation on itself and the trivial action on the affine
